\chapter{Le metaeuristiche, scolpire una soluzione}
\label{ch:metodo_metaeuristiche}

\epigraph{``Heuristics are mental shortcuts; meta-heuristics are
shortcuts about how to take shortcuts.''}{Christian e
Griffiths, \emph{Algorithms to Live By}~\parencite{christian2016}}

\lettrine[lines=3]{C}{P-SAT}, lo abbiamo visto nel capitolo
precedente, \`e bravissimo a trovare soluzioni fattibili, e
quando ottimizza una funzione obiettivo trova spesso anche
buone soluzioni in pochi minuti. Ma su istanze grandi e
complesse, dopo un certo punto, tende a fermarsi su un
\emph{plateau}: continua a frugare lo spazio delle soluzioni
senza migliorare significativamente quella corrente. Le
metaeuristiche sono una famiglia di tecniche complementari
che intervengono proprio quando questo accade: prendono una
soluzione gi\`a esistente e la migliorano iterativamente
esplorando il \emph{vicinato}, ovvero lo spazio delle
soluzioni ``simili'', in modo intelligente.

In $\pi$Tantum la Phase C della pipeline \`e dedicata
interamente alle metaeuristiche. Sono sei: LNS, ALNS,
Simulated Annealing, Tabu Search, Iterated Local Search e
Variable Neighborhood Search. Hanno parecchio in comune,
nel senso che sono tutte forme di ricerca locale, ma
differiscono profondamente nel modo in cui esplorano lo
spazio delle soluzioni e nel modo in cui accettano le mosse
nuove. Questo capitolo le presenta una alla volta, partendo
dalla loro radice comune, e chiude con una discussione di
quale convenga scegliere a seconda della situazione.

\section{La radice comune: la ricerca locale}

La ricerca locale, da cui tutte le metaeuristiche
discendono, \`e uno schema iterativo elementare. Si parte da
una soluzione iniziale, qualunque essa sia, ottenuta da un
algoritmo greedy o da una qualsiasi euristica costruttiva.
Si definisce un \emph{vicinato} della soluzione corrente,
ovvero l'insieme delle soluzioni ottenibili da quella
attuale con una singola mossa elementare; nel timetabling
una mossa tipica \`e lo scambio di due lezioni dello stesso
docente, oppure lo spostamento di una lezione in uno slot
libero. Si esplora il vicinato cercando una soluzione
migliore della corrente, e se la si trova vi ci si sposta;
poi si ripete dal nuovo punto. Quando il vicinato non
contiene pi\`u alcuna soluzione migliore, ci si ferma: si \`e
in un \emph{ottimo locale}.

Il problema della ricerca locale pura, prevedibilmente, \`e
proprio che si incastra nei minimi locali. Se l'ottimo
globale richiede di passare attraverso una soluzione
temporaneamente peggiore, la ricerca locale non la
attraverser\`a mai e si fermer\`a in una soluzione
sub-ottima. Le varie metaeuristiche sono altrettante
strategie per evitare o per uscire da questa trappola.

\section{Simulated Annealing: l'analogia con la metallurgia}

\begin{storiabox}
Nel 1983 Kirkpatrick, Gelatt e Vecchi pubblicarono su
\emph{Science} l'articolo \emph{Optimization by Simulated
Annealing}~\parencite{kirkpatrick1983}, importando
nell'ottimizzazione combinatoria un'idea presa direttamente
dalla metallurgia. Per ottenere una struttura cristallina
ordinata in un metallo non basta raffreddarlo bruscamente,
perch\'e in tal caso la materia rimane intrappolata in uno
stato disordinato; bisogna invece scaldarlo e raffreddarlo
lentamente, lasciandogli il tempo di trovare configurazioni
a bassa energia. La temperatura agisce come parametro di
``rumore termico'': alta temperatura permette movimenti
casuali che peggiorano la situazione, bassa temperatura
ammette soltanto movimenti migliorativi. L'idea \`e tradurre
questo processo fisico in un algoritmo di ricerca, che
all'inizio accetta movimenti peggiorativi con una certa
probabilit\`a e via via diventa pi\`u selettivo. L'articolo
metropolis del 1953~\parencite{metropolis1953} aveva fornito
il fondamento statistico, applicato in origine a simulazioni
di sistemi termodinamici.
\end{storiabox}

L'idea di simulated annealing si traduce in una regola
operativa molto semplice. Ad ogni iterazione si propone
una mossa: se essa migliora la funzione obiettivo, la si
accetta sempre. Se la peggiora di una quantit\`a $\Delta$,
la si accetta con probabilit\`a $e^{-\Delta / T}$, dove $T$
\`e la temperatura corrente. All'inizio $T$ \`e alta e
anche peggioramenti consistenti vengono accettati con buona
probabilit\`a; man mano che $T$ scende, soltanto i piccoli
peggioramenti passano; alla fine, quando $T$ \`e quasi zero,
il comportamento si avvicina a quello di una ricerca locale
greedy che accetta soltanto miglioramenti.

In $\pi$Tantum lo schema di raffreddamento di default \`e
geometrico, con $T_{k+1} = \alpha T_k$ e $\alpha = 0.95$. Si
parte da una temperatura iniziale tale che le mosse iniziali
abbiano una probabilit\`a di accettazione del settanta-ottanta
per cento, e si scende fino a $T \sim 0.01$, momento in cui
il sistema converge.

\section{Tabu Search: il taccuino delle mosse vietate}

\begin{storiabox}
Fred Glover propose Tabu Search nel 1986 con il celebre
articolo \emph{Future Paths for Integer Programming and
Links to Artificial Intelligence}~\parencite{glover1986},
oggi considerato uno dei lavori fondativi della moderna
metaeuristica. La parola \emph{tabu} viene dal polinesiano e
significa ``proibito'': l'idea \`e che, in ricerca locale,
dopo aver visitato una soluzione si vieti di tornarci per un
certo periodo, in modo da forzare la ricerca a esplorare
zone diverse dello spazio.
\end{storiabox}

Tabu Search mantiene una \emph{lista tabu} delle ultime
mosse fatte (per esempio ``ho spostato matematica della 1A
da luned\`i alla prima ora a marted\`i alla terza''). Le
mosse nella lista tabu non possono essere usate finch\'e non
escono dalla finestra: la lista \`e di lunghezza fissa, e
ogni nuova mossa spinge fuori la pi\`u vecchia. Esiste
inoltre un meccanismo detto \emph{aspiration criterion}: se
una mossa tabu produrrebbe la migliore soluzione mai vista
fino a quel momento, la si esegue lo stesso, perch\'e non
avrebbe senso rifiutare la migliore opzione disponibile per
il solo motivo che si trova nella lista. La lista tabu in
$\pi$Tantum \`e tipicamente di venti-quaranta mosse, valore
tarato sui benchmark dei profili di test.

\section{Iterated Local Search: perturbare e ricominciare}

L'idea di Iterated Local Search, sviluppata in modo
sistematico da Loureno, Martin e
Stuetzle~\parencite{lourenco2003}, \`e che una semplice
ricerca locale si ferma su un ottimo locale, ma rilanciandola
da una soluzione \emph{simile ma perturbata} si arriva
spesso a un ottimo locale diverso, eventualmente migliore.
Iterando il procedimento si esplorano molti bacini di
attrazione diversi senza dover ricominciare da zero ogni
volta.

L'algoritmo alterna due operazioni. La prima \`e una ricerca
locale tradizionale fino al raggiungimento di un ottimo
locale $x^*$. La seconda \`e una perturbazione: si applica a
$x^*$ un disturbo casuale di entit\`a controllata,
producendo una soluzione $x'$, di solito ottenuta facendo
cinque-dieci scambi casuali di lezioni. La ricerca locale
ricomincia da $x'$, e produce un nuovo ottimo locale $x^{**}$
che si confronta con $x^*$ tenendo il migliore dei due. ILS
\`e spesso il metaeuristico ``base'' che si combina con
altri (per esempio ILS+Tabu Search, ILS+Simulated Annealing)
per ottenere ibridi pi\`u potenti.

\section{Variable Neighborhood Search: cambiare obiettivo}

VNS, dovuto a Mladenovi\'c e Hansen del
1997~\parencite{mladenovic1997}, propone una strategia
diversa per uscire dai minimi locali: invece di perturbare
la soluzione corrente, si cambia il \emph{tipo} di vicinato
in cui si cerca. L'idea \`e che un ottimo locale rispetto a
un vicinato piccolo non \`e necessariamente un ottimo
locale rispetto a un vicinato pi\`u grande, e quindi
crescendo il raggio della ricerca si pu\`o uscire
dall'incastro.

In $\pi$Tantum sono definiti quattro vicinati di dimensione
crescente: il primo \`e l'1-swap, ovvero lo scambio di due
lezioni dello stesso docente; il secondo \`e il 2-swap, due
coppie indipendenti di lezioni scambiate insieme; il terzo
\`e il 3-chain, una rotazione di tre lezioni del tipo a$\to$b$\to$c$\to$a;
il quarto \`e il $k$-opt con $k$ fra quattro e sei, una
rotazione lunga che coinvolge molte lezioni
contemporaneamente. VNS parte sempre dal vicinato pi\`u
piccolo, e quando trova un miglioramento ricomincia dal
vicinato uno; quando un'intera passata da uno a quattro non
trova alcun miglioramento, l'algoritmo si ferma. La
soluzione finale \`e un ottimo locale rispetto a tutti e
quattro i vicinati.

\section{Large Neighborhood Search: distruggere e ricostruire}

LNS fu introdotto da Paul Shaw nel
1998~\parencite{shaw1998} per il problema del vehicle
routing, e da allora si \`e diffuso in molti altri domini
applicativi del scheduling combinatorio. L'idea \`e radicale
rispetto alle metaeuristiche tradizionali: invece di
esplorare un vicinato piccolo (uno o due swap), si
\emph{distrugge} una porzione significativa della soluzione
corrente (per esempio il venti per cento delle lezioni) e la
si \emph{ricostruisce ottimamente} con un solver esatto. Il
vicinato cos\`i definito \`e ``grande'', nel senso che
contiene tutte le soluzioni ottenibili ricostruendo la parte
distrutta in modo diverso, ma viene esplorato in modo
intelligente dal solver invece che con una scansione
elementare.

LNS alterna dunque due operatori. Il primo \`e
l'\emph{operator destroy}, che rimuove un sottoinsieme delle
variabili dalla soluzione corrente, lasciando le altre
fisse. Il secondo \`e l'\emph{operator repair}, che
re-ottimizza le variabili rimosse con tutti i vincoli del
problema attivi e con le altre variabili fisse come dato
del problema. In $\pi$Tantum l'operator repair \`e quasi
sempre CP-SAT con un budget di tempo di pochi secondi, e gli
operator destroy disponibili sono numerosi:
\texttt{random\_window} (rimuove un blocco di slot
contigui), \texttt{day\_cluster} (rimuove tutte le lezioni
di un intero giorno), \texttt{worst\_fit\_day} (rimuove il
giorno con il peggior valore obiettivo),
\texttt{teacher\_day} (rimuove tutte le lezioni di un
docente in un giorno), \texttt{classroom\_day} e
\texttt{single\_class\_week} (varianti analoghe). Gli
operator repair classici sono \texttt{cp\_sat\_window}
(re-risolve con CP-SAT), \texttt{greedy\_by\_soft} (ricostruisce
con un greedy che ottimizza l'obiettivo SOFT) e
\texttt{bfs\_fill\_back} (ricostruisce con una ricerca in
ampiezza guidata).

\section{Adaptive LNS: scolpire con scalpelli adattivi}

\begin{storiabox}
ALNS fu proposto da Ropke e Pisinger nel
2006~\parencite{ropke2006} estendendo LNS con un meccanismo
di apprendimento. Invece di scegliere a caso fra gli
operator destroy e repair disponibili, ALNS impara durante
l'esecuzione quali funzionano meglio sull'istanza corrente
e li seleziona pi\`u spesso. Il meccanismo si ispira agli
algoritmi di apprendimento per rinforzo: ad ogni iterazione
ogni operatore riceve un punteggio proporzionale al successo
della sua applicazione, e i punteggi vengono mediati nel
tempo con una decay esponenziale.
\end{storiabox}

L'analogia che aiuta a capire ALNS \`e quella dello
scultore davanti a un blocco di marmo: lo scultore non
usa sempre lo stesso scalpello, ma ne ha tanti, e sceglie
quello giusto a seconda della parte della statua su cui sta
lavorando. Allo stesso modo ALNS sceglie ad ogni iterazione
l'operator destroy e l'operator repair pi\`u promettenti
sulla base dei punteggi accumulati, calcolati come media
pesata esponenzialmente decrescente nel tempo. Il criterio
di accettazione delle mosse \`e tipicamente analogo a quello
di simulated annealing.

L'aggiornamento dei punteggi avviene in modo semplice. Dopo
ogni iterazione, l'operatore usato riceve un premio, di
ammontare $\sigma_1$, $\sigma_2$ o $\sigma_3$ a seconda del
risultato: $\sigma_1$ se la mossa ha migliorato globalmente
la migliore soluzione mai vista, $\sigma_2$ se ha migliorato
soltanto la soluzione corrente, $\sigma_3$ se \`e stata
accettata via meccanismo SA-like pur peggiorando. Il
punteggio dell'operatore $w_i$ viene aggiornato come
$w_i \leftarrow \rho w_i + (1 - \rho) \sigma$ con
$\rho = 0.95$. La selezione dell'operatore avviene poi via
roulette wheel pesata sui punteggi correnti.

\section{Quale scegliere}

\begin{ideabox}
Una regola pratica orientativa pu\`o essere la seguente. Se
si ha poco tempo a disposizione e si vuole un risultato
``decente'' in pochi minuti, conviene usare simulated
annealing con uno schema di raffreddamento veloce. Se si
vuole esplorare in modo robusto lo spazio, ILS o VNS sono
buone scelte. Se si ha molto tempo a disposizione e si vuole
un risultato vicino all'ottimo, ALNS \`e quasi sempre il
migliore perch\'e impara a fare le mosse giuste sull'istanza
corrente. Se la soluzione iniziale \`e molto strutturata e
servono mosse grandi, LNS con CP-SAT come repair offre un
buon compromesso fra ampiezza dell'esplorazione e qualit\`a
del risultato. In Phase C, $\pi$Tantum esegue per default
ALNS, e mette opzionalmente a disposizione le altre per
sperimentazione e confronto.
\end{ideabox}

\section{Un esempio concreto su otto lezioni}

\begin{esempiobox}
Riprendiamo l'esempio del cap.~\ref{ch:metodo_cpsat} con
tre docenti, due classi e otto lezioni. Una soluzione
fattibile \`e gi\`a stata trovata da CP-SAT con valore della
funzione obiettivo $f(x_0) = 5$ (somma pesata di ore isolate
dei docenti, per esempio).

Iterazione 1: simulated annealing prova uno scambio fra le
ore 1 e 2 della matematica della 1A. Il nuovo valore \`e
$f(x_1) = 4$. Migliora! Accetta sempre.

Iterazione 2: prova uno scambio fra l'italiano della 1A
(marted\`i prima) e la storia della 1A (luned\`i seconda).
Il nuovo valore \`e $f(x_2) = 6$. Peggiora di $\Delta = 2$.
Con $T = 1.0$ la probabilit\`a di accettazione \`e $P =
e^{-2/1.0} = 0.135$. Estrae un numero casuale $r = 0.4 > P$
e rifiuta la mossa. Resta su $x_1$.

Iterazione 3: prova lo scambio fra la prima ora di
matematica della 1A e quella della 1B. Il nuovo valore \`e
$f = 3$. Migliora, accetta.

Schema di raffreddamento: $T \leftarrow 0.95 T = 0.95$.

Cos\`i di seguito per migliaia di iterazioni. Dopo
cinquemila iterazioni tipiche $T < 0.01$ e il sistema
converge a un ottimo locale rispetto al vicinato 1-swap. Su
questa istanza piccola il valore finale \`e quasi sempre
$f = 2$, che si pu\`o dimostrare essere l'ottimo globale.
\end{esempiobox}

\section{Limiti e parametri da tarare}

\begin{avvertenzabox}
Le metaeuristiche non garantiscono l'ottimo. Su problemi
piccoli possono trovare soluzioni peggiori di CP-SAT esatto
in poche iterazioni; conviene quindi combinarle, applicando
prima CP-SAT esatto fino al raggiungimento del plateau e poi
le metaeuristiche per il raffinamento. I parametri da
tarare sono diversi: la temperatura iniziale di simulated
annealing, la lunghezza della lista tabu di tabu search, la
dimensione della perturbazione di ILS. $\pi$Tantum ha valori
di default ragionevoli per ciascun profilo di test, ma su
istanze atipiche conviene fare uno sweep di tuning. Le mosse
infine devono essere economicamente valutabili: $\pi$Tantum
implementa una \emph{delta-evaluation} della funzione
obiettivo, che evita di ricalcolare tutto da zero ad ogni
iterazione e che porta il costo per swap da $O(n)$ a $O(1)$.
\end{avvertenzabox}

\section{Le connessioni con il resto del sistema}

Le metaeuristiche non vivono isolate. CP-SAT
(cap.~\ref{ch:metodo_cpsat}) viene usato come solver di
\emph{repair} in LNS e in ALNS, perch\'e la sua capacit\`a
di re-ottimizzare velocemente sotto-vicinati piccoli
costituisce il motore della distruzione-ricostruzione. La
decomposizione spettrale (cap.~\ref{ch:metodo_spettrale})
viene applicata prima delle metaeuristiche, che lavorano poi
sopra la soluzione gi\`a parzialmente strutturata. L'editor
manuale dell'orario nell'interfaccia utente sfrutta gli
stessi schemi di mossa di simulated annealing: gli scambi
dell'utente sono mosse 1-swap del medesimo vicinato, e
l'editor mostra in tempo reale il delta della funzione
obiettivo esattamente come fa internamente l'algoritmo.

\begin{biblioparvabox}
Per chi voglia approfondire i contenuti di questo capitolo
in modo sistematico, segnalo alcuni riferimenti chiave.
L'articolo originale di simulated
annealing~\parencite{kirkpatrick1983} si legge in tre
pagine. L'articolo originale di tabu
search~\parencite{glover1986} \`e di Glover. L'articolo
originale di LNS~\parencite{shaw1998} \`e di Shaw, e
quello originale di ALNS~\parencite{ropke2006} \`e di
Ropke e Pisinger. Per una visione d'insieme, l'\emph{Handbook
of Metaheuristics}~\parencite{glover2003} edito da Glover e
Kochenberger \`e il riferimento collettivo. Van Hentenryck e
Michel, \emph{Constraint-Based Local
Search}~\parencite{vanhentenryck2005}, propongono una sintesi
elegante fra constraint programming e ricerca locale.
\end{biblioparvabox}
